\begin{lemma}[The next hierarchy edge uses an old reserved color]
Let a hierarchy be good and let $1\le j\le q$.  The first edge added in
$T_{n,j+1}-T_{n,j}$ has a color $\alpha$ for which there is an
$\eta_g\in D_{n,j}$ such that
\[
             \alpha\in\Gamma^{j-1}_{g}
             \quad\hbox{and}\quad
             \alpha\notin\Gamma^{j}_{g}.
\]
\end{lemma}
\begin{proof}
Write
\[
 A=T_{n,j}=H.\operatorname{stage}(i+1),\qquad
 B=T_{n,j+1}=H.\operatorname{stage}(i+2),
 \qquad j=i+1.
\]
Thus $i$ is an index below the number of dividers.  The nestedness and
strictness fields of \noderef{TashkinovHierarchy}, applied at the exact
index $i+1$, give an ordered initial segment $A\subset B$ and
$|A|<|B|$.  Consequently
\[
                  r:=|A|
\]
is a valid edge index of $B$, and $B.e(r)$ is exactly the first displayed
edge after $A$.  Iterating the same ordered initial-segment fields and
using the recorded terminal-stage equality gives $B\subseteq T$.

Let $n$ be the last series index, put
$T_n=\mathcal T.\operatorname{tree}(n)$, and denote the certified output
above it by $T_{n+1}$.  By \noderef{ExtendedTashkinovTree}, $T$ is an
ordered initial segment of $T_{n+1}$.  Separately, the base-stage field
of \noderef{TashkinovHierarchy} identifies its stage zero with $T_n$;
iterating nestedness therefore gives $T_n\subseteq A$.  Hence
\[
             T_n\subseteq A\subset B\subseteq T\subseteq T_{n+1}
\tag{1}
\]
as ordered tree-sequences, with equality of every edge identity in the
overlapping positions.  This last entrywise interpretation is precisely
\noderef{TreeSequenceInitialSegment}.

The series certificate supplied by \noderef{TashkinovSeries} has a
\noderef{TashkinovSeriesStep} at the positive terminal index $n$; it is
therefore one of the PE, SE, or RE constructors.  We record its
kind-dependent augmentation seed.
\begin{itemize}
\item In the PE case, the constructor says directly that $T_{n+1}$ is a
closure of the seed $T_n$ under the terminal coloring $\varphi_n=\varphi$.
By (1), this seed is an ordered initial segment of $A$.
\item In the SE case, let $f_n$ be the single connecting edge and let
$T_n^+$ be the constructor's displayed adjoined tree.  The certificate
$\operatorname{TreeSequenceAdjoinEdge}(T_n,f_n,T_n^+)$ in
\noderef{TreeSequenceAdjoinEdge} says exactly that
$|T_n^+|=|T_n|+1$, that $T_n\subset T_n^+$ in order, and that the new
last edge identity is $f_n$.  The SE constructor says that $T_{n+1}$ is
a closure of this exact seed $T_n+f_n$.
\item The RE case is identical at the level needed here: its singleton
connecting edge is the RE certificate's $f_n$, its displayed tree
$T_n^+$ satisfies the same exact
$\operatorname{TreeSequenceAdjoinEdge}(T_n,f_n,T_n^+)$ relation, and
the RE constructor makes $T_{n+1}$ a closure of $T_n^+$.
\end{itemize}
In the SE and RE cases, the strict first hierarchy inclusion from
\noderef{TashkinovHierarchy} gives $|T_n|<|T_{n,1}|\leq |A|$.
Thus $|T_n^+|=|T_n|+1\leq|A|$.  The closure chain makes $T_n^+$ an
ordered initial segment of $T_{n+1}$, while (1) does the same for $A$.
The entrywise equalities in \noderef{TreeSequenceInitialSegment}, first
for their common terminal tree and then restricted to the shorter
length, show that $T_n^+\subseteq A$.  Therefore, in all three cases,
the constructor's closure seed is an ordered initial segment of $A$.

We next justify that the exact edge $B.e(r)$ is a Tashkinov augmentation
edge from $A$.  Unfolding \noderef{TashkinovClosure} gives a
\noderef{TashkinovAugmentationSequence} from the appropriate seed to
$T_{n+1}$.  This is a finite reflexive-transitive chain of
\noderef{TashkinovAugmentation} steps.  Each such step increases the
length by exactly one and preserves every earlier ordered vertex and edge
entry.  Since the seed length is at most $|A|$ and, by (1),
$|A|<|B|\leq|T_{n+1}|$, the list of successive lengths contains $|A|$
and then $|A|+1$.  Truncate the chain at those two consecutive members,
say $Q$ and $Q^+$.  Both $Q$ and $A$ are ordered initial segments of
$T_{n+1}$ of length $|A|$, so the entrywise clause of
\noderef{TreeSequenceInitialSegment} identifies all of their vertices
and edges.  Similarly, $Q^+$ and the first $|A|+1$ entries of $B$ agree
with the same terminal sequence.  Hence the new edge in the augmentation
$Q\to Q^+$ is the individual edge identity
\[
                         e:=B.e(r).
\tag{2}
\]
By \noderef{TashkinovAugmentation}, there is therefore a color $\alpha$
such that $\varphi(e)=\alpha$ and
$\alpha\in\overline\varphi(A)$.

The same augmentation also makes $e$ a boundary edge of $A$, as we now
check rather than suppressing the incidence argument.  The new vertex at
position $r+1$ is outside $V(A)$: otherwise the injectivity clause of
\noderef{TreeSequence}, together with the entrywise vertex equalities in
\noderef{TreeSequenceInitialSegment}, would identify it with one of the
earlier vertices.  The attachment clause of \noderef{TreeSequence} says
that $e$ joins this new vertex to a vertex in an earlier position, hence
to a vertex of $A$.  By the finite vertex-set definition
\noderef{TreeSequenceVertexSet} and the exact edge-boundary definition
\noderef{BoundaryEdges}, $e\in\partial_G(V(A))$.  Combining this with
$\varphi(e)=\alpha$, the filtered definition
\noderef{BoundaryColorEdges} gives
\[
 e\in\partial_{\varphi,\alpha}(V(A)).
\tag{3}
\]

Set
\[
 C=\bigcup_{\eta\in D_{n,j}}\Gamma^{j-1}_{\eta}.
\]
Clause (v) of \noderef{GoodHierarchy}, at the exact divider index $i$,
says that $A$ is $C^{-}$-closed.  By
\noderef{TreeColorSetMinusClosed}, this means that $A$ is closed for
every color in $\overline\varphi(A)-C$.  If $\alpha\notin C$, the
already proved membership $\alpha\in\overline\varphi(A)$ places
$\alpha$ in that difference.  Unfolding \noderef{TreeColorSetClosed}
would then make the $\alpha$-colored boundary empty, contradicting the
specific member (3) and the definition in
\noderef{BoundaryColorEdges}.  Therefore $\alpha\in C$.  Unpacking this
finite union supplies a color $\eta_g$ with
\[
 \eta_g\in D_{n,j}
 \quad\hbox{and}\quad
 \alpha\in\Gamma^{j-1}_{g}.
\tag{4}
\]

It remains to exclude $\alpha$ from the new reserved pair.  Unfolding
\noderef{HierarchyStoredColorsAt}, the first membership in (4) says that
$\eta_g$ belongs to the terminal stored-color set but not to the tree
missing-color set of the starred stage at $j$.  Since $j=i+1>0$, the
positive-stage clause of \noderef{HierarchyStarredStage} identifies that
starred stage exactly with $A$.  Thus
$\eta_g\notin\overline\varphi(A)$.  By the union-over-vertices definition
\noderef{TreeMissingColors}, $\eta_g$ is not missing at any vertex of
$A$.

Let $p$ be the first missing position of $\eta_g$ in the terminal tree
$T$, as defined in \noderef{HierarchyFirstMissingPosition}.  We claim
\[
                             |A|<p.
\tag{5}
\]
If the finite set of missing positions is nonempty, $p$ is its least
member.  Were $p\leq |A|$, the ordered inclusion $A\subseteq T$ from
(1), unpacked with \noderef{TreeSequenceInitialSegment}, would identify
the vertex $T(p)$ with a vertex of $A$.  It would be missing $\eta_g$,
contrary to the preceding paragraph.  Hence (5) holds in this case.  If
there is no missing position, \noderef{HierarchyFirstMissingPosition}
sets $p$ equal to the last vertex position $|T|$.  Here
$|A|<|B|\leq|T|$ by strictness, terminal nestedness, and (1), so (5)
again holds.

Now (2), $r=|A|$, (5), and $|A|<|B|$ give
\[
 |A|\leq r<\min\{p,|B|\},\qquad \varphi(B.e(r))=\alpha.
\]
The exact filter in \noderef{HierarchySegmentColors}, with its starred
lower stage identified by \noderef{HierarchyStarredStage}, therefore
places $\alpha$ in the color set of the segment
$T_{n,j+1}(v_{\eta_g})-T_{n,j}$.  Finally clause (i), namely the
avoidance field of \noderef{GoodHierarchy} at $(j,\eta_g)$, makes
$\Gamma^j_g$ disjoint from this segment-color set.  Its contrapositive
at the established member $\alpha$ yields
$\alpha\notin\Gamma^j_g$.  Together with (2) and (4), this proves both
assertions with the required first edge identity.
\end{proof}
